% =====================================================================
%  Project Thống Kê — bài thuyết trình (Beamer, biên dịch bằng xelatex)
%    cd presentation && xelatex slides.tex && xelatex slides.tex
%  Hình: analysis/figs1.py, figs2.py, figs3.py  ->  figs/*.pdf
%  Số liệu trên slide: analysis/numbers1.json, numbers2.json, numbers3.json
%  Font: fonts/BeVietnamPro-*.ttf (đi kèm thư mục)
% =====================================================================
\documentclass[aspectratio=169,10pt,t]{beamer}
\usetheme{metropolis}
\metroset{sectionpage=none, numbering=none, progressbar=none, titleformat frame=regular}
\usepackage{fontspec}
\usepackage{polyglossia}
\setdefaultlanguage{vietnamese}
\setsansfont{BeVietnamPro}[Path=fonts/, Extension=.ttf, UprightFont=*-Regular, BoldFont=*-SemiBold,
  ItalicFont=*-Italic, BoldItalicFont=*-SemiBold]
\setmainfont{BeVietnamPro}[Path=fonts/, Extension=.ttf, UprightFont=*-Regular, BoldFont=*-SemiBold, ItalicFont=*-Italic]
\newfontfamily\heavy{BeVietnamPro}[Path=fonts/, Extension=.ttf, UprightFont=*-Bold]
\newfontfamily\medium{BeVietnamPro}[Path=fonts/, Extension=.ttf, UprightFont=*-Medium]
\setmonofont{Latin Modern Mono}
\usepackage{booktabs,tabularx,array,colortbl}
\usepackage[most]{tcolorbox}
\usepackage{tikz}
\usetikzlibrary{positioning,arrows.meta,calc}

% ---------------------------------------------------------------- bảng màu pastel
\definecolor{surf}{HTML}{FBFBFD}
\definecolor{ink}{HTML}{1E2233}
\definecolor{ink2}{HTML}{555B70}
\definecolor{muted}{HTML}{8A90A3}
\definecolor{hair}{HTML}{E6E8F0}
\definecolor{goldsoftx}{HTML}{FBF2DC}
% mỗi bước một màu: tone (nét), tint (nền), dark (chữ)
\definecolor{toneN}{HTML}{8A90A3}\definecolor{tintN}{HTML}{EEF0F5}\definecolor{darkN}{HTML}{454B61}
\definecolor{toneA}{HTML}{578ED4}\definecolor{tintA}{HTML}{E6EEFA}\definecolor{darkA}{HTML}{2E5A93}
\definecolor{toneB}{HTML}{EE9370}\definecolor{tintB}{HTML}{FCEBE3}\definecolor{darkB}{HTML}{9E4E2E}
\definecolor{toneC}{HTML}{279479}\definecolor{tintC}{HTML}{DFF2EB}\definecolor{darkC}{HTML}{1B6A56}
\definecolor{toneD}{HTML}{D7AC4A}\definecolor{tintD}{HTML}{FBF2DC}\definecolor{darkD}{HTML}{85631A}
\definecolor{toneE}{HTML}{8C6EBD}\definecolor{tintE}{HTML}{EFEAF8}\definecolor{darkE}{HTML}{5D4690}
\definecolor{toneF}{HTML}{D87DA2}\definecolor{tintF}{HTML}{FAE5EE}\definecolor{darkF}{HTML}{944565}
\definecolor{toneI}{HTML}{279479}\definecolor{tintI}{HTML}{E8F5F0}\definecolor{darkI}{HTML}{1B6A56}
\definecolor{ord1}{HTML}{BFD2F3}\definecolor{ord2}{HTML}{8DABE4}\definecolor{ord3}{HTML}{6380C8}\definecolor{ord4}{HTML}{41559E}
\colorlet{steptone}{toneN}\colorlet{steptint}{tintN}\colorlet{stepdark}{darkN}
\gdef\stepkicker{}
\newcommand{\setstep}[2]{\colorlet{steptone}{tone#1}\colorlet{steptint}{tint#1}\colorlet{stepdark}{dark#1}\gdef\stepkicker{#2}}

\setbeamercolor{normal text}{fg=ink,bg=surf}
\setbeamercolor{background canvas}{bg=surf}
\setbeamercolor{alerted text}{fg=stepdark}
\setbeamercolor{itemize item}{fg=steptone}
\setbeamercolor{enumerate item}{fg=stepdark}
\setbeamercolor{palette primary}{fg=darkE,bg=tintE}
\setbeamercolor{standout}{fg=darkE,bg=tintE}
\setbeamertemplate{itemize item}{\color{steptone}\small$\bullet$}
\setbeamersize{text margin left=0.6cm, text margin right=0.6cm}
\setlength{\leftmargini}{1.0em}
\arrayrulecolor{hair}
\renewcommand{\arraystretch}{1.3}
\renewcommand{\tabularxcolumn}[1]{>{\raggedright\arraybackslash}p{#1}}
\newcolumntype{L}[1]{>{\raggedright\arraybackslash}p{#1}}

% ---------------------------------------------------------------- pill, kicker
\newcommand{\ra}{$\rightarrow$}
\newcommand{\pill}[3]{\tikz[baseline=(p.base)]\node[fill=#1,text=#2,rounded corners=3pt,inner xsep=4.5pt,inner ysep=1.6pt,font=\scriptsize\bfseries](p){\vphantom{Ếg}#3};}
\newcommand{\vU}{\pill{tintC}{darkC}{Ủng hộ}}
\newcommand{\vM}{\pill{tintD}{darkD}{Một phần}}
\newcommand{\vB}{\pill{tintF}{darkF}{Bác bỏ}}
\newcommand{\hyp}[2]{{\color{ink2}Giả thuyết: #1}\hspace{0.5em}\csname v#2\endcsname}
\newcommand{\ask}[1]{{\color{ink2}Câu hỏi: #1}}

% ---------------------------------------------------------------- tiêu đề frame
% dòng 1: kicker (bước · mã) · dòng 2: headline · dòng 3: giả thuyết + kết quả
\makeatletter
\newlength\barw \setlength\barw{\dimexpr\paperwidth-1.2cm\relax}
\setbeamertemplate{frametitle}{%
  \nointerlineskip
  \hbox{{\color{steptone}\rule{\dimexpr\paperwidth*\insertframenumber/\inserttotalframenumber\relax}{2.2pt}}%
        {\color{hair}\rule{\dimexpr\paperwidth-\paperwidth*\insertframenumber/\inserttotalframenumber\relax}{2.2pt}}}%
  \nointerlineskip
  \hbox to \paperwidth{\hspace{0.6cm}\vbox to 38pt{\hsize=\barw\vskip 5pt
      \hbox to \barw{\fontsize{13.5}{15}\selectfont\bfseries\color{ink}%
        \ifx\framecode\@empty\else\raisebox{1pt}{\pill{steptone!28!white}{stepdark}{\framecode}}\hspace{0.4em}\fi
        \strut\insertframetitle\hfil}%
      \vskip 1pt
      \hbox to \barw{\fontsize{7.8}{9}\selectfont\strut\insertframesubtitle\hfil}%
      \vfil}\hss}%
  \vskip 4pt}
\gdef\framemeta{}\gdef\framecode{}
\newcommand{\code}[1]{\gdef\framecode{#1}}
\setbeamertemplate{footline}{%
  \hbox to \paperwidth{\hspace{0.6cm}\fontsize{6.4}{8}\selectfont\color{muted}\framemeta\hfill\insertframenumber\,/\,\inserttotalframenumber\hspace{0.6cm}}%
  \vskip 6pt\gdef\framemeta{}\gdef\framecode{}}
\makeatother
\newcommand{\meta}[3]{\gdef\framemeta{\textbf{Thước đo}\enspace #1\qquad\textbf{Biểu đồ}\enspace #2\qquad\textbf{Kiến thức}\enspace #3}}

% ---------------------------------------------------------------- thẻ nhận xét
\tcbset{card/.style={enhanced, colback=steptint, colframe=steptint, boxrule=0pt, arc=5pt,
  left=5pt, right=5pt, top=3pt, bottom=3pt, fontupper=\footnotesize\color{ink}, before upper={\raggedright\setlength{\parskip}{2pt}}}}
\newenvironment{cardrow}[1]{\begin{tcbraster}[raster columns=#1, raster equal height, raster column skip=6pt, raster row skip=5pt, raster before skip=0pt, raster after skip=0pt]}{\end{tcbraster}}
\newcommand{\card}[1]{\begin{tcolorbox}[card]#1\end{tcolorbox}}
\newcommand{\num}[1]{{\bfseries\color{stepdark}#1}}

% ---------------------------------------------------------------- khung hình rộng + thẻ
\newsavebox\bodybox
\newlength\figh \newlength\framereserve \setlength\framereserve{78pt}
% \figcards{hình}{số cột}{\card{..}\card{..}}
\newcommand{\figcards}[3]{%
  \sbox\bodybox{\parbox[t]{\textwidth}{\begin{cardrow}{#2}#3\end{cardrow}}}%
  \setlength\figh{\dimexpr\textheight-\framereserve-\ht\bodybox-\dp\bodybox\relax}%
  \typeout{FIGH #1 \the\figh\space BODY \the\ht\bodybox\space + \the\dp\bodybox}%
  {\centering\includegraphics[width=\textwidth,height=\figh,keepaspectratio]{figs/#1.pdf}\par}%
  \vspace{5pt}\usebox\bodybox\par\vfill}
\newcommand{\sidefig}[1]{\includegraphics[width=\linewidth,height=\dimexpr\textheight-\framereserve+6pt\relax,keepaspectratio]{figs/#1.pdf}}
% hình trái, cột thẻ phải
\newcommand{\sidecards}[3]{%
  \begin{columns}[T]
  \begin{column}{#2\textwidth}\centering\sidefig{#1}\end{column}
  \begin{column}{\dimexpr0.97\textwidth-#2\textwidth\relax}\begin{cardrow}{1}#3\end{cardrow}\end{column}
  \end{columns}}

\begin{document}
\typeout{TH \the\textheight}

% ================================================================= Tiêu đề
\begin{frame}[plain]
\begin{tikzpicture}[remember picture, overlay]
  \fill[tintE] ($(current page.north east)+(-6.1,0)$) rectangle (current page.south east);
  % motif: median thu nhập/người 2014 của 4 nhóm Q (triệu đồng)
  \begin{scope}[shift={($(current page.south east)+(-5.35,1.55)$)}]
    \foreach \i/\v/\c/\l in {0/9.18/ord1/{9,2}, 1/18.66/ord2/{18,7}, 2/29.56/ord3/{29,6}, 3/51.55/ord4/{51,6}}{
      \fill[\c, rounded corners=3pt] (\i*1.12,0) rectangle ++(0.82,{\v*0.085});
      \node[font=\scriptsize\bfseries, text=darkE, anchor=south] at ({\i*1.12+0.41},{\v*0.085+0.05}) {\l};
      \node[font=\scriptsize, text=ink2, anchor=north] at ({\i*1.12+0.41},-0.08) {Q\the\numexpr\i+1\relax};
    }
    \node[font=\fontsize{6.5}{8}\selectfont, text=ink2, anchor=north west, align=left] at (-0.05,-0.55) {Median thu nhập/người năm 2014\\(triệu đồng) theo 4 nhóm thu nhập};
  \end{scope}
\end{tikzpicture}
\vspace*{1.35cm}
\begin{minipage}{0.56\textwidth}
{\fontsize{8}{10}\selectfont\medium\color{darkE} VHLSS 2010 · 2012 · 2014\par}
\vspace{6pt}
{\fontsize{34}{38}\selectfont\heavy\color{ink} Project\\Thống Kê\par}
\vspace{10pt}
{\large\color{ink} Hộ giàu, hộ nghèo chi tiêu\\và tích lũy thế nào?\par}
\vspace{12pt}
{\small\color{ink2} Khảo sát Mức sống Hộ gia đình Việt Nam\par\vspace{2pt}28.196 hộ\quad·\quad 63 tỉnh\quad·\quad 3 đợt khảo sát\par}
\vspace{1.3cm}
{\small\color{muted} Thành viên nhóm: ……………………………………\par}
\end{minipage}
\end{frame}

% ================================================================= Mở đầu
\setstep{N}{Mở đầu}
\section{Mở đầu}
\begin{frame}\frametitle{Túi tiền hay nơi ở quyết định cách chi tiêu?}
\framesubtitle{\ask{một câu hỏi, ba phạm vi so sánh, ba câu hỏi lặp lại ở mỗi phạm vi}}
\vspace{4pt}
\begin{tcbraster}[raster columns=3, raster equal height, raster column skip=8pt]
\begin{tcolorbox}[card, colback=tintA]{\scriptsize\medium\color{darkA} PHẠM VI 1}\par\vspace{2pt}{\large\bfseries Toàn quốc}\par\vspace{4pt}{\color{ink2} So sánh 4 nhóm thu nhập: Q1 = 25\% hộ có thu nhập/người thấp nhất … Q4 = 25\% cao nhất.}\end{tcolorbox}
\begin{tcolorbox}[card, colback=tintB]{\scriptsize\medium\color{darkB} PHẠM VI 2}\par\vspace{2pt}{\large\bfseries Thành thị / Nông thôn}\par\vspace{4pt}{\color{ink2} Cùng một nhóm Q, sống ở thành thị hay nông thôn có khác không?}\end{tcolorbox}
\begin{tcolorbox}[card, colback=tintE]{\scriptsize\medium\color{darkE} PHẠM VI 3}\par\vspace{2pt}{\large\bfseries Nhóm tỉnh T1–T4}\par\vspace{4pt}{\color{ink2} Cùng một nhóm Q, sống ở nhóm tỉnh nghèo (T1) hay giàu (T4) có khác không?}\end{tcolorbox}
\end{tcbraster}
\vspace{12pt}
{\small\color{ink2} Ở mỗi phạm vi, nhóm hỏi lại 3 câu:\par\vspace{4pt}
\textbf{\color{ink}(a)} Q1 khác Q4 thế nào?\qquad \textbf{\color{ink}(b)} Cùng Q, các nơi khác nhau thế nào?\qquad \textbf{\color{ink}(c)} Từ 2010 đến 2014 thay đổi ra sao?\par}
\end{frame}

\begin{frame}\frametitle{Mạch trình bày: 6 bước nối tiếp nhau}
\framesubtitle{\ask{mỗi bước trả lời một câu hỏi và dẫn sang câu hỏi tiếp theo}}
\vspace{6pt}
\centering
\begin{tikzpicture}[
  st/.style={rounded corners=5pt, minimum height=2.2cm, text width=1.86cm, align=left, inner sep=5pt, font=\footnotesize, anchor=north west},
  arr/.style={-{Stealth[length=4pt]}, color=muted, line width=0.7pt}]
\foreach \i/\L/\q/\m in {0/A/{Ai giàu, ai nghèo, ở đâu?}/A1–A4, 1/B/{Bao nhiêu cho thiết yếu?}/B1–B4, 2/C/{Tiền đi vào đâu?}/C1–C5, 3/D/{Tích lũy tài sản gì?}/D1–D5, 4/E/{Thay đổi ra sao 2010–14?}/E1–E5, 5/F/{Túi tiền hay nơi ở?}/F1–F4}{
  \node[st, fill=tint\L] (s\i) at ({\i*2.5},0) {{\scriptsize\medium\color{dark\L} BƯỚC \the\numexpr\i+1\relax}\\[2pt]{\bfseries\color{ink}\q}\\[5pt]{\scriptsize\color{dark\L}\m}};
}
\foreach \i in {0,...,4}{\draw[arr] ($(s\i.east)+(0.02,0.2)$) -- ++(0.28,0);}
\node[rounded corners=5pt, fill=tintN, font=\footnotesize, inner sep=5pt, anchor=north west, text width=6.9cm] at (0,-2.6) {{\medium\color{darkN} Trước đó:} {\color{ink} dữ liệu \ra\ làm sạch \ra\ chia nhóm}};
\node[rounded corners=5pt, fill=tintI, font=\footnotesize, inner sep=5pt, anchor=north west, text width=6.9cm] at (7.6,-2.6) {{\medium\color{darkI} Sau đó:} {\color{ink} chuyện bên lề \ra\ kết luận \ra\ hạn chế}};
\end{tikzpicture}
\par\vfill
\raggedright{\small\color{ink2} Mỗi slide phân tích: tiêu đề là \textbf{\color{ink}kết luận}; dòng dưới là \textbf{\color{ink}giả thuyết} và kết quả \vU\ \vM\ \vB; chân trang ghi thước đo, loại biểu đồ, kiến thức môn học.\par}
\vspace{4pt}
\end{frame}

% ================================================================= Dữ liệu
\setstep{N}{Dữ liệu và làm sạch}
\section{Dữ liệu và làm sạch}
\begin{frame}\frametitle{28.196 hộ, 3 đợt khảo sát, 63 tỉnh}
\framesubtitle{\ask{dữ liệu có đủ lớn để so sánh theo nhóm không?}}
\meta{tần số, tần suất}{pie chart, grouped bar}{C1 – dữ liệu thứ cấp, mẫu và tổng thể}
\figcards{f01_sample}{2}{%
\card{Dữ liệu \num{thứ cấp} của Tổng cục Thống kê: chính thức, mẫu lớn, nhưng không hỏi thêm được (không biết ai hút thuốc, ai có thẻ BHYT).}
\card{Mỗi đợt là một \num{mẫu khác nhau}, không có mã hộ, không có trọng số: so sánh theo nhóm, không theo dõi từng hộ.}}
\end{frame}

\begin{frame}\frametitle{30 biến: 13 khoản chi và 10 loại tài sản}
\framesubtitle{\ask{mỗi biến thuộc loại nào và đo cái gì?}}
\meta{—}{bảng}{C1 – thang đo, loại biến; C2 – tứ phân vị}
\footnotesize
\begin{tabularx}{\textwidth}{@{}L{2.7cm}X L{2.9cm}@{}}
\rowcolor{tintN}\textbf{Nhóm} & \textbf{Biến} & \textbf{Loại biến (C1)}\\
Phân loại & \texttt{province} (63 tỉnh), \texttt{area} (thành thị/nông thôn), \texttt{Quart} (1–4), \texttt{Year} & định danh, nhị phân, thứ bậc\\
Thu nhập & \texttt{size} (số người), \texttt{income}, \texttt{Inc\_head} = \texttt{income} / \texttt{size} & định lượng\\
Chi tiêu trong năm & \texttt{rice nrice tobaco nfood edu med hins aid water elec garb fuel fhouse} & định lượng liên tục\\
Tài sản đang có & \texttt{house oto moto bicycle tele tv emachine hmachine bed tab} & định lượng liên tục\\
\bottomrule
\end{tabularx}
\normalsize
\vspace{8pt}
\begin{cardrow}{3}
\card{\texttt{Quart}: tứ phân vị của \texttt{Inc\_head} trên toàn mẫu \num{từng năm}; mỗi nhóm luôn đúng 25\% số hộ.}
\card{Tài sản là \num{giá trị đồ đang có}, không phải tiền mua trong năm (94\% hộ có máy gia dụng).}
\card{Đơn vị gốc: nghìn đồng/năm. Trên slide, thu nhập ghi theo triệu đồng/người/năm.}
\end{cardrow}
\end{frame}

\begin{frame}\frametitle{Thu nhập lệch phải rất mạnh}
\framesubtitle{\ask{nên dùng mean hay median để mô tả thu nhập?}}
\meta{mean, median, skewness, quartile, IQR}{histogram}{C2 – vị trí, hình dạng, quy tắc outlier}
\figcards{f02_income}{2}{%
\card{Năm 2014: mean \num{29,98} triệu nhưng median chỉ \num{23,61} triệu đồng/người. Gộp 3 năm: skewness 33,7; 5,4\% hộ là outlier (trên Q3 + 1,5·IQR).}
\card{Trên thang log, phân phối gần đối xứng: các biểu đồ thu nhập về sau dùng \num{median} và \num{trục log}.}}
\end{frame}

\begin{frame}\frametitle{Làm sạch: không xóa hộ nào, điền 0 có gắn cờ}
\framesubtitle{\ask{ô trống nghĩa là “không có” hay “không biết”?}}
\meta{tần số, tỷ lệ}{bảng}{C3 – phương sai mẫu; C4 – ước lượng chệch}
\footnotesize
\begin{tabularx}{\textwidth}{@{}L{2.1cm}L{2.3cm}X L{3.4cm}@{}}
\rowcolor{tintN}\textbf{Biến} & \textbf{Số ô thiếu} & \textbf{Phát hiện} & \textbf{Xử lý}\\
\texttt{fuel} & 817 (813 ở 2014) & 2014 không hộ nào ghi 0; ô trống tập trung ở hộ nghèo nông thôn & điền 0 + cờ \texttt{fuel\_na}\\
9 cột tài sản & 284 hộ thiếu cả 9 & 89\% ở nông thôn, phần lớn hộ Q1; chỉ 2 hộ khai cả 9 cột bằng 0 & điền 0 + cờ \texttt{asset\_na}\\
\texttt{house}, \texttt{fhouse} & 23 hộ & rải rác, không rõ nguyên nhân & \texttt{house}: bỏ khi phân tích nhà; \texttt{fhouse} = 0\\
\texttt{nfood} & 1 hộ & — & bỏ khỏi phân tích chi tiêu\\
Giá trị âm & 61 ô / 36 hộ & $-1$, $-2$…: tài sản không thể âm & coi là ô thiếu\\
\bottomrule
\end{tabularx}
\normalsize
\vspace{8pt}
\begin{cardrow}{2}
\card{\textbf{Vì sao điền 0?} Số hộ ghi nhiên liệu = 0: 40 (2010), 199 (2012), \num{0} (2014), nhưng 2014 có 813 ô trống: năm 2014 để trống thay cho 0.}
\card{Tài sản: chỉ 2 hộ khai cả 9 cột bằng 0, trong khi hộ Q1 có 1, 2, 3 món là 5,8\%, 8,9\%, 11,3\%, nên phải có khoảng 3\% không có món nào: khớp \num{3,3\%} hộ Q1 bị trống.}
\end{cardrow}
\end{frame}

\begin{frame}\frametitle{Xóa dòng hay điền median/mean đều làm hộ nghèo “khá” lên}
\framesubtitle{\ask{sao không xóa dòng thiếu, hoặc điền median/mean cho đỡ lệch?}}
\meta{mean, median, phương sai}{bar chart}{C2 – mean và median; C4 – ước lượng chệch}
\sidecards{f03b_missing}{0.44}{%
\card{\textbf{Xóa dòng?} Ô trống dồn vào hộ nghèo (2014: Q1 \num{21,5\%}, Q4 1,4\%). Xóa thì nhiên liệu TB Q1 tăng \num{27\%}, khoảng cách TV Q4–Q1 hẹp lại (19,0 \ra\ 16,6 điểm).}
\card{\textbf{Điền median/mean?} Hộ Q1 có khai chỉ chi 2,4 triệu, còn median cả nước 4,32, mean 5,69 (lệch phải): Q1 tăng \num{40\%} và \num{53\%}. Điền mean còn làm phương sai nhỏ đi: $S^2_{\text{sau}} = \frac{n_{\text{có}}-1}{n-1}S^2_{\text{có}}$.}
\card{\textbf{Khi nào được xóa?} Khi ô thiếu ít và rải rác (23 hộ thiếu giá nhà). Bỏ hộ \texttt{fuel\_na}: Q1 thiếu tiền thiết yếu 2014 là 33,7\% thay vì 32,5\%, kết luận giữ nguyên.}}
\end{frame}

% ================================================================= Chia nhóm
\setstep{N}{Chia nhóm}
\section{Chia nhóm}
\begin{frame}\frametitle{Nhóm tỉnh T1–T4 theo hạng thu nhập}
\framesubtitle{\ask{làm sao so tỉnh giữa các năm mà không cần quy đổi giá?}}
\meta{median, xếp hạng, tứ phân vị}{bar chart}{C2 – quantile}
\figcards{f04_tiers}{2}{%
\card{Median \texttt{Inc\_head} của tỉnh mỗi năm \ra\ xếp hạng \ra\ hạng trung bình 3 năm \ra\ \num{4 nhóm} 16/16/15/16 tỉnh.}
\card{T1: miền núi phía Bắc, Bắc Trung Bộ. T4: Hà Nội, TP.HCM, Bình Dương, Đà Nẵng, Bắc Ninh… Ô nhỏ nhất có \num{39 hộ}.}}
\end{frame}

\begin{frame}\frametitle{Phân loại chi tiêu (Maslow) và tài sản (công dụng)}
\framesubtitle{\ask{khoản chi nào là “phải chi”? Tài sản là đồ đang có, nên chỉ chia theo công dụng}}
\meta{tỷ lệ}{bảng}{phân loại trước khi phân tích}
\begin{columns}[T]
\begin{column}{0.56\textwidth}
\footnotesize
\begin{tabularx}{\linewidth}{@{}L{3.3cm}L{1.1cm}X@{}}
\rowcolor{tintN}\textbf{Khoản chi} & \textbf{Tầng} & \textbf{Nhóm}\\
\texttt{rice nrice water elec fuel} & 1 & Thiết yếu\\
\texttt{garb med aid hins fhouse} & 2 & Thiết yếu\\
\texttt{nfood} (quần áo, đi lại…) & 3–4 & Không thiết yếu\\
\texttt{edu} & 4–5 & Không thiết yếu\\
\texttt{tobaco} & — & Không thiết yếu\\
\bottomrule
\end{tabularx}
\end{column}
\begin{column}{0.41\textwidth}
\footnotesize
\begin{tabularx}{\linewidth}{@{}L{2.1cm}X@{}}
\rowcolor{tintN}\textbf{Nhóm tài sản} & \textbf{Biến}\\
Nhà & \texttt{house}\\
Phương tiện & \texttt{oto moto bicycle}\\
Công nghệ & \texttt{tele tv emachine}\\
Đồ gia dụng & \texttt{hmachine}\\
Nội thất & \texttt{bed tab}\\
\bottomrule
\end{tabularx}
\end{column}
\end{columns}
\normalsize
\vspace{8pt}
\begin{cardrow}{3}
\card{Kiểm tra bằng số liệu: tỷ trọng \texttt{nfood} tăng \num{20\% \ra\ 26\%} từ Q1 tới Q4, đúng kiểu nhu cầu bậc cao.}
\card{Xếp \texttt{edu} vào thiết yếu thì tỷ trọng thiết yếu cao hơn 4–6 điểm nhưng vẫn giảm theo Q: kết luận không đổi.}
\card{Nhà chiếm \num{94–95\%} giá trị tài sản ở mọi Q, nên được phân tích riêng (D1).}
\end{cardrow}
\end{frame}

\begin{frame}\frametitle{Một khung phân tích cho cả 3 phạm vi}
\framesubtitle{\ask{so sánh thế nào cho công bằng?}}
\meta{—}{—}{C2, C3}
\vspace{2pt}
\begin{cardrow}{3}
\card{{\medium\color{darkN}(a) Q1 khác Q4?}\par Trong cùng phạm vi, trục x là Q1 … Q4.}
\card{{\medium\color{darkN}(b) Cùng Q, khác nơi?}\par Thành thị với nông thôn; T1 với T4 ở cùng một Q.}
\card{{\medium\color{darkN}(c) 2010 \ra\ 2014?}\par Từng phạm vi qua 3 năm (Bước 5).}
\end{cardrow}
\vspace{10pt}
\begin{itemize}\small
\setlength{\itemsep}{4pt}
\item Thước đo đủ các nhóm của C2: vị trí (mean, median, mode, sum, quantile), phân tán (range, variance, SD, CV, IQR), hình dạng (skewness), quan hệ (covariance, correlation); cùng tỷ lệ hộ (C3).
\item Ô dưới 30 hộ không kết luận. So sánh giữa các năm ưu tiên tỷ trọng, tỷ lệ hộ, thứ hạng, vì tiền ghi theo giá từng năm.
\item Mỗi giả thuyết được ghi kết quả: \vU\ \vM\ \vB.
\end{itemize}
\end{frame}

% ================================================================= Bước 1
\setstep{A}{Bước 1 · Ai giàu, ai nghèo, ở đâu}
\section{Bước 1 · Ai giàu, ai nghèo, ở đâu}
\begin{frame}\code{A1}\frametitle{Q4 có thu nhập gấp 5,6 lần Q1}
\framesubtitle{\hyp{ở mọi phạm vi, thu nhập điển hình của Q4 gấp khoảng 5 lần Q1}{U}}
\meta{mean, median, quantile q(0,1), q(0,9), tỷ số}{boxplot, bar chart}{C2}
\figcards{f05_A1}{2}{%
\card{2014: median Q1 \num{9,18} triệu, Q4 \num{51,55} triệu. Ở Q4 mean (62,35) lớn hơn hẳn median: nhóm giàu có đuôi rất dài.}
\card{q(0,9)/q(0,1): toàn quốc \num{6,9} lần; T1 \num{7,7}; thành thị chỉ 4,7. Hộ Q1 ở T1 nghèo nhất: median 8,3 triệu (T4: 10,0).}}
\end{frame}

\begin{frame}\code{A2}\frametitle{Hộ giàu ở thành thị, hộ nghèo ở nhóm tỉnh T1}
\framesubtitle{\hyp{hộ giàu tập trung ở thành thị và tỉnh giàu, hộ nghèo ở nông thôn và tỉnh nghèo}{U}}
\meta{tần số, tần suất theo hàng}{100\% stacked bar}{C2 – bảng chéo}
\sidecards{f06_A2}{0.55}{%
\card{Thành thị: \num{46\%} hộ thuộc Q4, chỉ 7\% thuộc Q1. Nông thôn: 32\% Q1, 16\% Q4.}
\card{T1: \num{48\%} Q1. T4: \num{41\%} Q4. Cơ cấu gần như không đổi qua 3 năm.}
\card{Vì vậy khi so thành thị với nông thôn phải so \num{trong cùng Q}, nếu không sẽ lẫn hiệu ứng “thành thị giàu hơn”.}}
\end{frame}

\begin{frame}\code{A3}\frametitle{Tỉnh nghèo có chênh lệch nội bộ lớn nhất}
\framesubtitle{\hyp{nơi giàu hơn (thành thị, T4) có chênh lệch thu nhập nội bộ lớn hơn}{B}}
\meta{SD, CV, correlation}{line chart, scatter plot}{C2 – phân tán, tương quan}
\figcards{f07w_A3}{2}{%
\card{CV (bỏ 1\% cao nhất): T1 cao nhất và \num{không giảm} (0,81 \ra\ 0,80); T4 giảm rõ 0,74 \ra\ 0,63, thành thị 0,70 \ra\ 0,61.}
\card{63 tỉnh: tỉnh càng giàu, CV càng thấp; r = $-0{,}53$ (2010) \ra\ \num{$-0{,}74$} (2014). Cao nhất Sơn La (1,14).}}
\end{frame}

\begin{frame}\code{A4}\frametitle{10\% hộ giàu nhất nhận 26\% tổng thu nhập}
\framesubtitle{\hyp{một nhóm nhỏ hộ giàu nhận phần lớn thu nhập, và khoảng cách giảm dần qua các năm}{U}}
\meta{sum, quantile (thập phân vị), tỷ lệ}{bar chart, line chart}{C2 – quantile, sum}
\figcards{f35_A4}{2}{%
\card{2014: 10\% hộ có thu nhập/người cao nhất (8,6\% dân số) nhận \num{26,3\%} tổng thu nhập; 40\% hộ thấp nhất (\num{42\% dân số}) chỉ nhận \num{17,4\%}.}
\card{10\% giàu nhất: 29,6 \ra\ 26,9 \ra\ 26,3\%; 40\% nghèo nhất: 17,2 \ra\ 17,4\%. Tỷ số giữa hai phần: 1,72 \ra\ \num{1,51 lần}. Thành thị đều hơn nông thôn.}}
\end{frame}

% ================================================================= Bước 2
\setstep{B}{Bước 2 · Tiền cho nhu cầu thiết yếu}
\section{Bước 2 · Tiền cho nhu cầu thiết yếu}
\begin{frame}\code{B1}\frametitle{Càng giàu, tỷ trọng thiết yếu càng giảm}
\framesubtitle{\hyp{hộ Q1 dành khoảng 3/4 chi tiêu cho thiết yếu, tỷ trọng giảm dần tới Q4}{U}}
\meta{sum, mean, median}{100\% stacked bar}{C2 – tần suất tương đối}
\sidecards{f08_B1a}{0.44}{%
\card{Tỷ trọng từ \num{sum} của cả nhóm: 74,4\% (Q1) \ra\ 72,2 \ra\ 70,6 \ra\ \num{67,3\%} (Q4).}
\card{Kiểm tra bằng tỷ trọng từng hộ: mean 75,9 \ra\ 69,0\%; median 76,6 \ra\ 69,9\%. Ba cách, một xu hướng.}
\card{Thu nhập Q4 gấp gần 6 lần Q1 mà cơ cấu \num{chỉ đổi 7 điểm}: khoản thiết yếu lớn nhất (thực phẩm ngoài gạo) không giảm theo thu nhập (B4).}}
\end{frame}

\begin{frame}\code{B1}\frametitle{Cùng Q, nơi ở gần như không đổi cơ cấu chi}
\framesubtitle{\hyp{ở cùng một Q, hộ nông thôn dành tỷ trọng thiết yếu cao hơn hộ thành thị}{M}}
\meta{sum (tỷ trọng)}{line chart theo Q}{C2}
\figcards{f09_B1b}{2}{%
\card{Thành thị và nông thôn chênh không quá \num{0,5 điểm} ở Q1–Q3; chỉ ở Q4 chênh 2,5 điểm (66,3\% so với 68,8\%).}
\card{Giữa các nhóm tỉnh: hộ Q1 ở T1 dành 76,0\% cho thiết yếu, ở T4 72,4\%. Thu nhập vẫn là yếu tố chính.}}
\end{frame}

\begin{frame}\code{B2}\frametitle{Hộ Q1 dùng 82\% thu nhập cho thiết yếu}
\framesubtitle{\hyp{hộ Q1 dùng phần lớn thu nhập cho thiết yếu, hộ Q4 dùng chưa tới một nửa}{U}}
\meta{median, quantile}{line chart, bar chart}{C2}
\figcards{f10t_B2}{2}{%
\card{Median 2014: Q1 \num{81,6\%} … Q4 \num{38,6\%}; giảm đều ở \num{mọi Q} từ 2010 (Q1: 88,6 \ra\ 81,6\%). Q3 của nhóm Q1 là 111\%.}
\card{Hộ Q1 ở T4 chịu gánh nặng cao nhất (\num{89,8\%}), có thể vì giá cả ở nhóm tỉnh giàu cao hơn (xem F2).}}
\end{frame}

\begin{frame}\code{B4}\frametitle{Gạo: từ 20\% xuống 6\% chi thiết yếu}
\framesubtitle{\hyp{thực phẩm chiếm phần lớn chi thiết yếu ở mọi Q, phần của gạo giảm dần}{U}}
\meta{sum, tỷ trọng}{pie chart}{C2 – tần suất tương đối}
\figcards{f11_B4pies}{1}{%
\card{Thiết yếu: thực phẩm ngoài gạo \num{57\% \ra\ 63\%}, gạo \num{20\% \ra\ 6\%} (Q1 \ra\ Q4). Qua thời gian, gạo trong chi thiết yếu của Q1 giảm đều 23 \ra\ 20 \ra\ \num{18\%}. Không thiết yếu: ngoài ăn uống 76–79\%, giáo dục 16–20\%, thuốc lá 5\% \ra\ 3\%.}}
\end{frame}

\begin{frame}\code{B4}\frametitle{Khoản chi thứ ba đổi từ gạo sang nhiên liệu}
\framesubtitle{\hyp{khoản chi lớn nhất của mỗi nhóm Q có thể khác nhau}{M}}
\meta{sum, tỷ trọng, xếp hạng}{horizontal bar}{C2}
\figcards{f12_B4top}{2}{%
\card{Hạng 1 luôn là thực phẩm ngoài gạo (\num{\textasciitilde 43\%} tổng chi), không giảm theo thu nhập: vì vậy tỷ trọng thiết yếu \num{chỉ giảm nhẹ}. Hạng 2 là ngoài ăn uống, 20\% \ra\ 26\%.}
\card{Hạng 3: gạo ở Q1–Q2 \ra\ nhiên liệu ở Q3–Q4. Ở nông thôn, T1 và T3, gạo giữ hạng 3 tới tận Q3.}}
\end{frame}

% ================================================================= Bước 3
\setstep{C}{Bước 3 · Tiền đi vào đâu}
\section{Bước 3 · Tiền đi vào đâu}
\begin{frame}\code{C1}\frametitle{Luật Engel ở Việt Nam là “luật gạo”}
\framesubtitle{\hyp{càng giàu, tỷ trọng thực phẩm càng giảm, chủ yếu do gạo}{U}}
\meta{sum (tỷ trọng), mean, median, SD, CV}{100\% stacked bar, grouped bar}{C2}
\figcards{f13_C1}{2}{%
\card{Thực phẩm 57,8\% (Q1) \ra\ 47,0\% (Q4): gần như toàn bộ phần giảm là gạo (\num{15,1\% \ra\ 4,3\%}); thực phẩm khác đứng yên quanh 43\%.}
\card{Tiền gạo/người gần như bằng nhau ở mọi Q (median \num{1,31–1,36} triệu, CV khoảng 0,4); nông thôn cao hơn thành thị.}}
\end{frame}

\begin{frame}\code{C2}\frametitle{Điện theo túi tiền, nước và rác theo nơi ở}
\framesubtitle{\hyp{tiền điện tăng theo Q; trả tiền nước, phí rác hay không phụ thuộc nơi ở}{U}}
\meta{median, tỷ lệ hộ}{line chart, grouped bar}{C2; C3 – tỷ lệ mẫu}
\figcards{f14t_C2}{2}{%
\card{Điện/người 2014: Q1 200 nghìn, Q4 \num{960 nghìn}. Từ 2010, Q1 tăng \num{×2,4}, Q4 ×1,9: khoảng cách Q4/Q1 hẹp lại 5,9 \ra\ 4,8 lần.}
\card{Cùng là Q1, tỷ lệ hộ trả tiền nước ở thành thị gấp \num{3,6 lần} nông thôn (43\% so với 12\%). Phí rác: thành thị 32 \ra\ 87\%, nông thôn 10 \ra\ 35\%.}}
\end{frame}

\begin{frame}\code{C3}\frametitle{Y tế ngược chiều giữa thành thị và nông thôn}
\framesubtitle{\hyp{hộ nghèo có tỷ trọng y tế cao hơn; hộ giàu, hộ thành thị chi nhiều hơn cho giáo dục}{M}}
\meta{sum (tỷ trọng), tỷ lệ hộ}{grouped bar}{C2}
\figcards{f15_C3}{2}{%
\card{Y tế: thành thị 7,2\% \ra\ 5,9\%, nông thôn \num{6,4\% \ra\ 7,9\%}. Hộ Q1 ở T4 chi 8,5\%, ở T1 chỉ 4,9\%.}
\card{Giáo dục: thành thị 4,7 \ra\ 6,5\%; nông thôn lên rồi xuống. Tính trên hộ có chi, giáo dục/người của Q4 gấp \num{5,9 lần} Q1; khoảng cách hẹp dần 6,5 \ra\ 6,1 \ra\ 5,4 lần.}}
\end{frame}

\begin{frame}\code{C4}\frametitle{Hộ giàu chi thêm ngoài ăn uống, bớt thuốc lá}
\framesubtitle{\hyp{ngoài ăn uống tăng theo Q; thuốc lá nặng hơn với hộ nghèo, hộ nông thôn}{U}}
\meta{sum (tỷ trọng), tỷ lệ hộ}{bar chart, grouped bar, line chart}{C2}
\figcards{f16t_C4}{2}{%
\card{Ngoài ăn uống \num{20,1\% \ra\ 26,0\%} (Q1 \ra\ Q4). Thuốc lá 1,31\% \ra\ 0,92\% tổng chi; \% hộ có chi thuốc lá \num{giảm ở mọi Q} (Q1: 71 \ra\ 65\%).}
\card{Bất ngờ: T2 có tỷ trọng thuốc lá cao nhất (1,62\%), T1 thấp nhất (0,71\%); có thể do thuốc lào tự trồng ở miền núi (chưa kiểm chứng).}}
\end{frame}

\begin{frame}\code{C5}\frametitle{Khoản chi nào đi theo thu nhập?}
\framesubtitle{\hyp{chi ngoài ăn uống, điện tăng mạnh theo thu nhập; gạo thì không}{U}}
\meta{hệ số tương quan r, tỷ lệ hộ}{heatmap tương quan}{C2; C3 – n $\geq$ 30}
\sidecards{f36s_C5}{0.6}{%
\card{r(ln khoản chi/người, ln thu nhập/người), hộ có chi, 2014: \num{ngoài ăn uống 0,71}, thực phẩm khác 0,70, \num{điện 0,67} đi sát thu nhập nhất, ở mọi phạm vi.}
\card{\num{Gạo: r = 0,01}: tiền gạo mỗi người không phụ thuộc thu nhập (“luật gạo”). Thuốc lá (0,15), BHYT (0,11) cũng gần như không đổi. Thứ tự giữ nguyên qua 2010, 2012, 2014.}}
\end{frame}

% ================================================================= Bước 4
\setstep{D}{Bước 4 · Tích lũy tài sản}
\section{Bước 4 · Tích lũy tài sản}
\begin{frame}\code{D1}\frametitle{Cùng Q, nhà thành thị đắt gấp 2–3 lần}
\framesubtitle{\hyp{nhà thành thị đáng giá hơn; hộ nghèo thành thị có nhà đắt hơn hộ giàu nông thôn}{M}}
\meta{mean, median, quantile, IQR}{grouped bar, boxplot}{C2}
\figcards{f17_D1}{2}{%
\card{Nhà Q1 thành thị 245 triệu, gấp 2,5 lần Q1 nông thôn, nhưng \num{thấp hơn} Q4 nông thôn (400 triệu). Thành thị: mean 1,26 tỷ, median 700 triệu.}
\card{Nhà/thu nhập: thành thị \num{5,7 năm}, nông thôn 3,1. Hộ Q1 thành thị: \num{3,9 \ra\ 6,4 năm} (2010 \ra\ 2014), hộ Q4 thành thị: 6,5 \ra\ 5,8. Nhà xa dần tầm với hộ nghèo.}}
\end{frame}

\begin{frame}\code{D2}\frametitle{Hộ nghèo ở đâu cũng có khoảng 1 triệu đồ công nghệ}
\framesubtitle{\hyp{hộ nghèo ở nơi giàu hơn (thành thị, T4) có đồ công nghệ đáng giá hơn}{B}}
\meta{median}{grouped bar, line chart}{C2}
\figcards{f18t_D2}{2}{%
\card{Median đồ công nghệ của hộ Q1: \num{0,95–1,1 triệu} ở mọi phạm vi. Khác biệt chỉ có ở hộ giàu: Q4 thành thị 5,45 triệu, nông thôn 3,4.}
\card{Qua thời gian: Q1 \num{đứng yên} ở 1,0 triệu, Q4 tăng 3,2 \ra\ \num{4,3 triệu}. Tỷ lệ có điện thoại thu hẹp (E1) nhưng khoảng cách về giá trị lại nới rộng.}}
\end{frame}

\begin{frame}\code{D3}\frametitle{Đồ công nghệ đi cùng thu nhập: r = 0,49}
\framesubtitle{\hyp{giá trị đồ công nghệ tăng cùng thu nhập}{U}}
\meta{covariance, correlation, mean, median}{scatter plot}{C2; C1 – loại biến}
\sidecards{f19_D3}{0.55}{%
\card{cov(ln thu nhập, ln công nghệ) = 0,387; chia cho 0,728 × 1,075 được \num{r = 0,49}. Mọi phạm vi: 0,43–0,49.}
\card{Đồ công nghệ luôn bằng \num{2,4–2,9\%} thu nhập năm ở mọi Q (mean 1,37 \ra\ 7,06 triệu).}
\card{Nhưng \textbf{có} điện thoại, TV không còn là dấu hiệu giàu: từ Q2 trở lên 88–97\% hộ đều có.}}
\end{frame}

\begin{frame}\code{D4}\frametitle{Bậc thang xe: xe đạp \ra\ xe máy \ra\ ô tô}
\framesubtitle{\hyp{khi giàu lên, phần phương tiện tăng, phần nội thất và đồ công nghệ giảm}{U}}
\meta{sum (cơ cấu), tỷ lệ hộ}{100\% stacked bar, line chart}{C2; C3}
\figcards{f20_D4}{2}{%
\card{Phương tiện chiếm 54,5\% giá trị đồ dùng ở Q1 và \num{61,9\%} ở Q4; nội thất 16,0 \ra\ 10,6\%, công nghệ 12,9 \ra\ 9,6\%.}
\card{Xe máy 65 \ra\ \num{92\%} hộ (Q1 \ra\ Q4); xe đạp lên rồi xuống: hộ giàu bỏ xe đạp. Hộ Q1 có xe máy tăng \num{53 \ra\ 67\%} (2010 \ra\ 2014).}}
\end{frame}

\begin{frame}\code{D5}\frametitle{Nhà và đồ gia dụng phản ánh giàu nghèo rõ nhất}
\framesubtitle{\hyp{nhà và phương tiện (xe) là dấu hiệu giàu rõ nhất}{M}}
\meta{hệ số tương quan r, tỷ lệ hộ}{heatmap tương quan}{C2; C3 – n $\geq$ 30}
\figcards{f21_D5}{2}{%
\card{r(ln giá trị nhóm, ln thu nhập/người), trên hộ có nhóm đó: \num{nhà 0,57}, \num{đồ gia dụng 0,55}, công nghệ 0,49, nội thất 0,41, phương tiện 0,39. Thứ tự gần như giữ nguyên ở mọi phạm vi và cả 3 năm.}
\card{Phương tiện \num{thấp nhất}: xe máy gần nhà nào cũng có, ô tô chỉ 2\% hộ. Cộng mọi đồ dùng lại (0,56) thì phản ánh giàu nghèo gần bằng giá trị nhà.}}
\end{frame}

% ================================================================= Bước 5
\setstep{E}{Bước 5 · Xu hướng 2010–2014}
\section{Bước 5 · Xu hướng 2010–2014}
\begin{frame}\code{E1}\frametitle{Nhóm giữa tăng nhanh nhất; hộ nghèo bớt gạo}
\framesubtitle{\hyp{tỷ trọng thiết yếu của hộ nghèo giảm dần từ 2010 đến 2014}{M}}
\meta{median, sum (tỷ trọng), \% tăng}{line chart}{C2}
\figcards{f22_E1a}{2}{%
\card{Median \texttt{Inc\_head} 2010 \ra\ 2014: Q1 +70,5\%, \num{Q2 +80,6\%}, Q3 +76,1\%, Q4 +66,7\% (giá từng năm: chỉ để so giữa các nhóm).}
\card{Thiết yếu của Q1: 75,5 \ra\ 76,8 \ra\ \num{71,7\%}; giảm rõ năm 2014 vì gạo giảm, ngoài ăn uống tăng. Bỏ các hộ trống nhiên liệu vẫn 71,0\%.}}
\end{frame}

\begin{frame}\code{E1}\frametitle{Mọi phạm vi cùng đi một hướng, tỉnh nghèo đi nhanh hơn}
\framesubtitle{\hyp{thu nhập tăng, chi thiết yếu và gạo giảm, đồ hiện đại phổ cập ở mọi phạm vi}{U}}
\meta{median, sum (tỷ trọng), tỷ lệ hộ}{line chart (lưới 3 phạm vi × 4 chỉ số)}{C2}
\figcards{f38_E1grid}{1}{%
\card{Cả 3 phạm vi cùng hướng: thu nhập tăng; \num{\% thiết yếu}, \num{tỷ trọng gạo} giảm; điện thoại tăng. Đường của Q1, nông thôn, T1 dốc hơn: nhóm nghèo đang bắt kịp.}}
\end{frame}

\begin{frame}\code{E2}\frametitle{Nông thôn đang bắt kịp thành thị}
\framesubtitle{\hyp{nông thôn tăng thu nhập nhanh hơn và bắt kịp dần về dịch vụ công}{U}}
\meta{median, tỷ số, tỷ lệ hộ}{line chart}{C2}
\figcards{f24_E2}{2}{%
\card{Median \texttt{Inc\_head}: nông thôn \num{+79\%}, thành thị +64\%; tỷ số thành thị/nông thôn 1,92 \ra\ 1,81 \ra\ \num{1,76}.}
\card{Phí rác nông thôn 17 \ra\ 30\%: khoảng cách với thành thị giảm từ 54 xuống 47 điểm. Điện thoại nông thôn 72 \ra\ 87\%.}}
\end{frame}

\begin{frame}\code{E3}\frametitle{Các tỉnh gần nhau hơn về tỷ lệ, xa hơn về số tiền}
\framesubtitle{\hyp{nhóm tỉnh nghèo tăng nhanh hơn T4, nên khoảng cách giữa các tỉnh thu hẹp}{M}}
\meta{median, tỷ số, CV, range}{line chart}{C2 – phân tán}
\figcards{f25_E3}{2}{%
\card{T1 +81\%, T2 +79\%, T3 +87\%, T4 \num{+69\%}: nhóm giàu nhất tăng chậm nhất; T4/T1 từ 2,33 xuống 2,17 lần.}
\card{CV của 63 median tỉnh giảm \num{0,361 \ra\ 0,320}, nhưng range tăng 24,3 \ra\ 35,1 triệu: range phụ thuộc mức của dữ liệu, CV thì không.}}
\end{frame}

\begin{frame}\code{E4}\frametitle{Tỉnh nghèo tăng thu nhập nhanh nhất}
\framesubtitle{\hyp{có vài tỉnh tăng và leo hạng vượt trội; các tỉnh đứng đầu vẫn giữ vị trí}{U}}
\meta{median, mean, SD, min, max, correlation}{horizontal bar}{C2}
\sidecards{f26a_E4}{0.5}{%
\card{63 mức tăng: mean 78\%, median 79,5\%, SD 21 điểm; thấp nhất Bắc Kạn 15\%, cao nhất \num{Điện Biên 123\%}.}
\card{Top 5: Điện Biên, Hòa Bình, Thái Nguyên, Hà Giang, Thanh Hóa; \num{4/5} thuộc T1.}
\card{TP.HCM +50\%, Bình Dương +51\% nhưng vẫn hạng 1–2; TP.HCM gấp 5,1 lần Lai Châu.}
\card{r giữa mức tăng 2010–12 và 2012–14 là \num{$-0{,}42$}: “tăng mạnh rồi chững lại” phần lớn là nhiễu do mẫu nhỏ (xem F4).}}
\end{frame}

\begin{frame}\code{E4}\frametitle{Thái Nguyên leo 15 bậc, Vĩnh Long tụt 27 bậc}
\framesubtitle{\hyp{có vài tỉnh tăng và leo hạng vượt trội; các tỉnh đứng đầu vẫn giữ vị trí}{U}}
\meta{median, xếp hạng}{line chart (thứ hạng qua 3 năm)}{C2}
\sidecards{f26b_E4}{0.52}{%
\card{Leo: \num{Thái Nguyên 34 \ra\ 19}, Nam Định 23 \ra\ 11, Khánh Hòa 30 \ra\ 18.}
\card{Tụt: \num{Vĩnh Long 14 \ra\ 41}, Kon Tum 28 \ra\ 55 \ra\ 48, Hưng Yên 10 \ra\ 27.}
\card{Leo đều qua 2 giai đoạn: Khánh Hòa, Thái Nguyên. Kon Tum tụt rồi hồi lại: mẫu nhỏ nên thứ hạng dao động.}}
\end{frame}

\begin{frame}\code{E5}\frametitle{Cái gì thay đổi nhiều nhất từ 2010 đến 2014?}
\framesubtitle{\hyp{chi cho gạo giảm, chi cho điện và nhu cầu bậc cao tăng; đồ hiện đại phổ cập nhanh}{U}}
\meta{sum (tỷ trọng), tỷ lệ hộ, thay đổi điểm \%}{dot-arrow chart (2010 \ra\ 2014)}{C2; C3}
\figcards{f39_E5}{2}{%
\card{Tổng chi: \num{gạo $-1{,}9$ điểm}, thực phẩm khác +1,8, ngoài ăn uống +1,4, điện +0,6. *Nhiên liệu 2014 có ô trống.}
\card{Đồ dùng: \num{điện thoại +14 điểm}, xe máy +11, TV +11; máy giải trí đứng yên (bỏ giường, bàn ghế: 2010 ghi khác).}}
\end{frame}

% ================================================================= Bước 6
\setstep{F}{Bước 6 · Túi tiền hay nơi ở}
\section{Bước 6 · Túi tiền hay nơi ở}
\begin{frame}\code{F1}\frametitle{Túi tiền quyết định chi tiêu, nơi ở quyết định dịch vụ}
\framesubtitle{\hyp{cơ cấu chi do thu nhập quyết định; nước máy, thu gom rác do nơi ở quyết định}{U}}
\meta{hiệu của tỷ trọng và tỷ lệ giữa các nhóm (2014)}{horizontal grouped bar}{C2}
\figcards{f27_F1}{2}{%
\card{Mỗi thanh là trung bình hai hiệu: Q4 − Q1 (cùng khu vực) và thành thị − nông thôn (cùng Q). Cơ cấu chi, BHYT chênh theo \num{thu nhập}; nước máy theo \num{nơi ở}; phí rác theo cả hai.}
\card{Đồ công nghệ: +3,4 triệu theo thu nhập, so với +1,0 triệu theo nơi ở. Qua thời gian, cả 3 phạm vi đi cùng hướng.}}
\end{frame}

\begin{frame}\code{F2}\frametitle{Hơn 1/3 hộ Q1 không đủ tiền cho thiết yếu}
\framesubtitle{\hyp{cùng là hộ Q1, sống ở nơi giàu hơn (thành thị, T4) khó trang trải thiết yếu hơn}{M}}
\meta{tỷ lệ hộ, sai số chuẩn}{bar chart, line chart}{C2; C3 – tỷ lệ mẫu}
\figcards{f33_F2}{2}{%
\card{“Không đủ sống” = chi thiết yếu > thu nhập: \num{36,1\%} hộ Q1 (Q4: 2,3\%). \num{T4: 41\%}, cao nhất; T2: 32,5\%; thành thị ≈ nông thôn (34,7\% và 36,2\%).}
\card{T4 hơn T2 8,5 điểm, lớn hơn nhiều sai số chuẩn $\sqrt{p(1-p)/n}$ (T4: 1,7; T2: 1,1 điểm). 2010 \ra\ 2014: T1 41,5 \ra\ \num{30,3\%}, T4 38,4 \ra\ 40,5\%.}}
\end{frame}

\begin{frame}\code{F3}\frametitle{Ai để dành được? Hộ giàu nông thôn nhiều nhất}
\framesubtitle{\hyp{hộ giàu để dành tỷ lệ lớn hơn; hộ thành thị để dành nhiều hơn hộ nông thôn}{M}}
\meta{median, sum, tỷ lệ}{grouped bar, line chart}{C2}
\figcards{f37_F4}{2}{%
\card{Để dành = (thu nhập − tổng chi)/thu nhập. Median Q1 \num{$-14\%$} … Q4 \num{39\%}. Cùng Q4: nông thôn \num{46\%}, thành thị 33\%.}
\card{2010 \ra\ 2014 mọi nhóm để dành nhiều hơn (Q4: 36 \ra\ 43\%). Tổng chi chỉ gồm 13 khoản: đây là mức tối đa.}}
\end{frame}

\begin{frame}\code{F4}\frametitle{Hai kết quả gây ấn tượng nhất phần lớn là ảo ảnh}
\framesubtitle{\hyp{tỉnh tăng mạnh thì sau đó chững lại; hộ nghèo chi vượt thu nhiều hơn hẳn hộ giàu}{B}}
\meta{covariance, correlation, tỷ lệ}{scatter plot, grouped bar}{C2 – covariance; C3 – phân phối mẫu}
\figcards{f40_F4}{2}{%
\card{\num{Tỉnh:} nếu chỉ có nhiễu do mẫu (\textasciitilde 140 hộ/tỉnh), r giữa hai giai đoạn phải là \num{$-0{,}50$}; thực tế $-0{,}42$. Chia đôi ngẫu nhiên hộ mỗi tỉnh: hai nửa gần như không khớp (r = 0,14).}
\card{\num{Hộ:} xếp theo thu nhập, \% chi vượt thu giảm 64 \ra\ 9\%; xếp theo chi tiêu thì phẳng (28–34\%). Đó là \num{hồi quy về trung bình}.}}
\end{frame}

% ================================================================= Chuyện bên lề
\setstep{I}{Chuyện bên lề}
\section{Chuyện bên lề}
\begin{frame}\code{I1–I2}\frametitle{Hộ tự đóng BHYT lại chi y tế nhiều hơn}
\framesubtitle{\ask{ai tự mua bảo hiểm y tế, và mua rồi có bớt chi viện phí không?}}
\meta{tỷ lệ hộ, median, correlation}{line chart, grouped bar}{C2; C3}
\figcards{f28t_I12}{2}{%
\card{Tự đóng BHYT tăng ở Q2–Q4 (Q2: 44 \ra\ \num{54\%}), Q1 \num{đứng yên} khoảng 23\%: nhà nước đóng cho hộ nghèo (Luật BHYT 2008, Điều 12).}
\card{Hộ có BHYT chi y tế/thu nhập cao hơn (Q1: \num{3,7\%} so với 1,8\%): người hay ốm mới mua. \num{1/7 hộ} chi y tế > 10\% tổng chi, ở mọi Q.}}
\end{frame}

\begin{frame}\code{I3}\frametitle{Hút thuốc không đi cùng viện phí trong năm}
\framesubtitle{\ask{hộ hút thuốc nhiều có tốn viện phí nhiều hơn không?}}
\meta{correlation}{scatter plot}{C2}
\sidecards{f29_I3}{0.55}{%
\card{r trong từng nhóm Q: \num{0,09; $-0{,}01$; 0,00; $-0{,}01$}. Tác hại của thuốc lá đến sau nhiều năm, nên không đủ căn cứ nói “hút thuốc thì tốn viện phí”.}
\card{Nhưng hộ Q1 chi thuốc lá bằng \num{31\%} số tiền chi giáo dục (Q4: 16\%); 30\% hộ Q1 có cả hai khoản chi thuốc lá nhiều hơn giáo dục.}}
\end{frame}

\begin{frame}\code{I5–I6}\frametitle{Phí rác đo hạ tầng; HN, TP.HCM nhà đắt gấp 3,7}
\framesubtitle{\ask{dịch vụ đô thị đi theo điều gì, và hai siêu đô thị khác biệt đến đâu?}}
\meta{median, tỷ lệ, correlation}{scatter plot, bar chart}{C2}
\figcards{f30_I56}{2}{%
\card{Giữa 63 tỉnh: r(tỷ lệ hộ thành thị, \% trả phí rác) = 0,53; r(ln thu nhập tỉnh, \% phí rác) = \num{0,65}. 2010 \ra\ 2014 phí rác tăng ở cả 4 nhóm tỉnh (T1: 17 \ra\ 26\%).}
\card{Hà Nội + TP.HCM: median 40,4 triệu (T1–T3: 20,2), 57\% hộ thuộc Q4, 90\% trả phí rác; căn nhà bằng \num{11 năm} thu nhập (T1–T3: 3 năm).}}
\end{frame}

\begin{frame}\code{I7–I8}\frametitle{Hộ giàu ít người hơn}
\framesubtitle{\ask{vì sao hộ Q4 ít chi giáo dục hơn, và thu nhập đi cùng giáo dục đến đâu?}}
\meta{mode, mean, tỷ lệ, correlation}{bar chart, bubble chart}{C2 – tương quan}
\figcards{f31_I78}{2}{%
\card{Mode số người/hộ là 4 ở mọi Q (31\% số hộ), nhưng mean giảm: Q1 4,24; Q4 \num{3,48}. \num{44\%} hộ từ 7 người thuộc Q1. Hộ 2 người chỉ 14\% có chi giáo dục.}
\card{Tương quan thu nhập – chi giáo dục: \num{0,74–0,77} giữa 63 tỉnh, chỉ 0,47–0,52 giữa các hộ. Không suy từ tỉnh xuống từng hộ.}}
\end{frame}

\begin{frame}\code{I9–I11}\frametitle{Ba nghịch lý về tiền}
\framesubtitle{\ask{những con số nào đi ngược trực giác?}}
\meta{tỷ lệ, sum}{line chart, horizontal bar, bar chart}{C2}
\figcards{f32t_I9}{2}{%
\card{\num{64\%} hộ Q1 chi nhiều hơn thu (Q4: 9\%), một phần do cách chia nhóm (F4). Sum 2014: Q1 có 8,4\% tổng thu nhập nhưng 13,9\% tổng chi.}
\card{\num{77\%} hộ Q1 có TV, nhiều hơn điện thoại (67\%) và xe máy (65\%). Hộ Q4 ở \num{T1} có ô tô nhiều nhất (7,1\%).}}
\end{frame}

% ================================================================= Kết luận
\setstep{N}{Kết luận}
\section{Kết luận}
\begin{frame}\frametitle{Sáu thông điệp chính}
\framesubtitle{\ask{nếu chỉ nhớ một điều: túi tiền quyết định cách chi, nơi ở quyết định dịch vụ và chi phí sống}}
\vspace{2pt}
\begin{tcbraster}[raster columns=3, raster equal height, raster column skip=7pt, raster row skip=7pt]
\begin{tcolorbox}[card, colback=tintF]{\scriptsize\medium\color{darkF} F1}\par{\bfseries Túi tiền quyết định cách chi tiêu,} nơi ở quyết định dịch vụ công.\end{tcolorbox}
\begin{tcolorbox}[card, colback=tintB]{\scriptsize\medium\color{darkB} B2 · F2}\par{\bfseries Hộ Q1 dùng 82\% thu nhập cho thiết yếu;} hơn 1/3 hộ Q1 không đủ tiền cho thiết yếu, nhiều nhất ở T4.\end{tcolorbox}
\begin{tcolorbox}[card, colback=tintC]{\scriptsize\medium\color{darkC} C1 · B4}\par{\bfseries Luật Engel là “luật gạo”:} gạo từ 15\% xuống 4\% tổng chi, thực phẩm khác đứng yên quanh 43\%.\end{tcolorbox}
\begin{tcolorbox}[card, colback=tintD]{\scriptsize\medium\color{darkD} D3 · D5}\par{\bfseries Nhà và đồ gia dụng phản ánh giàu nghèo rõ nhất} (r 0,57 và 0,55); có TV hay điện thoại không còn là dấu hiệu giàu.\end{tcolorbox}
\begin{tcolorbox}[card, colback=tintA]{\scriptsize\medium\color{darkA} A3 · A4}\par{\bfseries 10\% hộ giàu nhất nhận 26\% tổng thu nhập;} tỉnh càng nghèo, chênh lệch nội bộ càng lớn (r = $-0{,}74$).\end{tcolorbox}
\begin{tcolorbox}[card, colback=tintE]{\scriptsize\medium\color{darkE} E2 · E3}\par{\bfseries Nông thôn và T1–T3 tăng nhanh hơn:} cách biệt tương đối thu hẹp, cách biệt tuyệt đối vẫn nới rộng.\end{tcolorbox}
\end{tcbraster}
\vspace{5pt}
\begin{tcolorbox}[card, colback=goldsoftx, fontupper=\small\color{ink}]
{\bfseries\color{darkD} Kết luận đắt giá:} khác biệt giữa các \textbf{nhóm lớn} (Q, khu vực, nhóm tỉnh) là thật; \textbf{thứ hạng từng tỉnh} và “\textbf{hộ nghèo chi vượt thu}” phần lớn là ảo ảnh của mẫu nhỏ và cách chia nhóm (F4).
\end{tcolorbox}
\end{frame}

\begin{frame}\frametitle{Hạn chế}
\framesubtitle{\ask{kết quả này nên đọc với những lưu ý nào?}}
\vspace{2pt}
\begin{itemize}\small
\setlength{\itemsep}{4pt}
\item Không có trọng số mẫu: kết quả là của mẫu VHLSS, không phải số chính thức của cả nước.
\item Mỗi năm là một mẫu khác nhau: chỉ so sánh theo nhóm, không theo dõi được từng hộ.
\item Tiền theo giá từng năm: \% tăng chỉ để so giữa các nhóm; ưu tiên tỷ trọng, tỷ lệ và thứ hạng.
\item Chi vượt thu trong một năm có thể được bù bằng tiết kiệm, vay mượn, trợ cấp; hộ rơi vào Q1 vì thu nhập năm đó thấp bất thường vẫn chi như thường. Vì vậy tỷ lệ “không đủ sống” (F2) có thể cao hơn thực tế; tổng chi chỉ gồm 13 khoản nên tỷ lệ để dành (F3) là mức tối đa.
\item Giá trị tài sản do hộ tự ước tính, theo số tròn: 90\% giá nhà là bội số của 10 triệu, mode 200 triệu. \texttt{hins} là tiền tự đóng, không cho biết hộ có thẻ BHYT hay không.
\item Tương quan không phải nhân quả; tương quan cấp tỉnh khác cấp hộ; mỗi tỉnh chỉ khoảng 140 hộ/năm nên số liệu từng tỉnh dao động mạnh (F4).
\item Không làm kiểm định giả thuyết (ngoài phạm vi môn); mọi kết luận dựa trên thống kê mô tả.
\end{itemize}
\end{frame}

\begin{frame}\frametitle{Kiến thức môn học đã vận dụng}
\framesubtitle{\ask{mỗi chương được dùng ở đâu trong bài?}}
\footnotesize
\begin{tabularx}{\textwidth}{@{}L{3.0cm}X L{4.4cm}@{}}
\rowcolor{tintN}\textbf{Chương} & \textbf{Nội dung} & \textbf{Ở đâu}\\
C1 Giới thiệu & dữ liệu thứ cấp; mẫu và tổng thể; thang đo, loại biến (không tính r với biến thứ bậc) & Dữ liệu, D3\\
C2 Thống kê mô tả & bảng tần số, bảng chéo; pie, bar, stacked bar, line, histogram, boxplot, scatter, bubble & mọi slide phân tích\\
 & mean, median, mode, sum, quantile (tứ, thập phân vị), range, variance, SD, CV, IQR, skewness & A1–A4, B1–B4, C1, E3, E4, F3, I7\\
 & covariance, hệ số tương quan (kể cả heatmap) & A3, C5, D3, D5, E4, F4, I2, I3, I5, I8\\
C3 Phân phối mẫu & tỷ lệ mẫu; sai số chuẩn $\sqrt{p(1-p)/n}$; quy tắc n $\geq$ 30; phương sai mẫu & C2, D4, F2, F4, chia nhóm, làm sạch\\
C4 Ước lượng & ước lượng chệch (điền mean làm phương sai nhỏ đi); median bền hơn mean khi có outlier & làm sạch, thu nhập\\
\bottomrule
\end{tabularx}
\end{frame}

\begin{frame}[standout,c]
\centering{\heavy\Large Cảm ơn thầy cô và các bạn đã lắng nghe!}\par\vspace{10pt}
{\normalsize\normalfont Hỏi đáp}
\end{frame}

\end{document}
